\clearpage\pagestyle{otMain}\setlength{\otBlocEspacement}{2.5mm}\fontsize{11}{15}\selectfont\corrigesChapitre{19}{La proportionnalité en actes}
\begin{blocLivre}{cor-19-1}\begin{corrige}{1}{Proportionnel ou pas ?}\demarche{Tester la phrase « fois plus » sur chaque situation.}
Baguettes : $\res{\text{oui}}$. Âge et pointure : $\res{\text{non}}$, doubler l’âge ne double pas la pointure. Allure régulière : $\res{\text{oui}}$, 2 fois plus de temps, 2 fois plus de distance.
\end{corrige}\end{blocLivre}
\begin{blocLivre}{cor-19-2}\begin{corrige}{2}{Fois plus}\demarche{Chercher le lien multiplicatif entre 4 et 12.}
12 croissants, 3 fois plus : $5\times3=\res{15}$ euros.
\end{corrige}\end{blocLivre}
\begin{blocLivre}{cor-19-3}\begin{corrige}{3}{Fois moins}\demarche{Relier chaque quantité à la donnée de départ.}
2 L, 3 fois moins : $30:3=\res{10}$ m². 18 L, 3 fois plus : $30\times3=\res{90}$ m².
\end{corrige}\end{blocLivre}
\begin{blocLivre}{cor-19-4}\begin{corrige}{4}{Additionner deux situations}\demarche{Assembler les deux achats connus.}
2 stylos, 2 fois moins : $\res3$ euros. Et $6=4+2$ stylos : $6+3=\res9$ euros.
\end{corrige}\end{blocLivre}
\begin{blocLivre}{cor-19-5}\begin{corrige}{5}{La recette au départ}\demarche{Combiner « fois plus », « fois moins » et sommes.}
18 crêpes : 3 fois plus, $150\times3=\res{450}$ g. 9 crêpes $=6+3$ : $150+75=\res{225}$ g, avec 3 crêpes à 75 g (2 fois moins que 6).
\end{corrige}\end{blocLivre}
\begin{blocLivre}{cor-19-6}\begin{corrige}{6}{Le débit de la pompe}\demarche{Découper 200 L en briques de 80 et 40.}
12 min, 3 fois plus : $80\times3=\res{240}$ L. Pour 200 L : $200=80+80+40$, soit $4+4+2=\res{10}$ min, avec 40 L en 2 min (2 fois moins).
\end{corrige}\end{blocLivre}
\begin{blocLivre}{cor-19-7}\begin{corrige}{7}{Le prix au kilogramme}\demarche{Revenir à l’unité quand elle simplifie tout.}
1 kg, 3 fois moins : $5{,}40:3=\res{1{,}80}$ euro. 7 kg : $7\times1{,}80=\res{12{,}60}$ euros ; ou $6$ kg $+1$ kg $=10{,}80+1{,}80$.
\end{corrige}\end{blocLivre}
\begin{blocLivre}{cor-19-8}\begin{corrige}{8}{L’échelle du plan}\demarche{Compter les paquets de 50 mètres.}
7 cm : $50\times7=\res{350}$ m. Réciproquement, 350 m, c’est 7 fois 50 m : $\res7$ cm sur le plan.
\end{corrige}\end{blocLivre}
\begin{blocLivre}{cor-19-9}\begin{corrige}{9}{Le piège du forfait}\demarche{Calculer les deux montants complets.}
3 activités : $8+6=14$ euros ; 6 activités : $8+12=20$ euros. Or $14\times2=28\ne20$ : $\res{\text{non proportionnel}}$, la part fixe casse le « fois plus ».
\end{corrige}\end{blocLivre}
\begin{blocLivre}{cor-19-10}\begin{corrige}{10}{Les deux imprimantes}\demarche{Comparer sur une même durée multiple commune.}
En 12 min : A, 4 fois plus, imprime $120\times4=480$ pages ; B, 3 fois plus, $150\times3=450$ pages. $\res{\text{A est la plus rapide}}$.
\end{corrige}\end{blocLivre}
\begin{blocLivre}{cor-19-11}\begin{corrige}{11}{La course régulière}\demarche{Découper la durée en briques simples.}
2 h : $24\times2=48$ km ; 30 min, 2 fois moins qu’une heure : 12 km. Total : $48+12=\res{60}$ km.
\end{corrige}\end{blocLivre}
\begin{blocLivre}{cor-19-12}\begin{corrige}{12}{Les photocopies de l’école}\demarche{Assembler des paquets de 100 et de 50.}
$250=100+100+50$ : $4+4+2=\res{10}$ euros, avec 50 copies à 2 euros (2 fois moins que 100).
\end{corrige}\end{blocLivre}
\begin{blocLivre}{cor-19-13}\begin{corrige}{13}{La maquette}\demarche{Multiplier les longueurs par le même facteur.}
12 m, 3 fois plus que 4 m : $8\times3=\res{24}$ cm. 2 m, 2 fois moins : $8:2=\res4$ cm. La réduction uniforme est une proportionnalité.
\end{corrige}\end{blocLivre}
\begin{blocLivre}{cor-19-14}\begin{corrige}{14}{Comparer les lots}\demarche{Ramener les deux offres au même nombre de cahiers.}
Deux lots A donnent 10 cahiers pour $7\times2=14$ euros, contre $13{,}50$ euros en B : $\res{\text{B est plus avantageux}}$ de $0{,}50$ euro.
\end{corrige}\end{blocLivre}
\begin{blocLivre}{cor-19-15}\begin{corrige}{15}{Proportionnel, le périmètre ?}\demarche{Tester chaque grandeur séparément.}
Périmètres : $4\times2=8$ et $4\times6=24=8\times3$ : le périmètre $\res{\text{triple}}$, il est proportionnel au côté. Aires : $4$ et $36=4\times9$ : l’aire est multipliée par $\res9$ : non proportionnelle au côté.
\end{corrige}\end{blocLivre}
